\chapter{Background and Literature Review}
\label{ch:review}
\vspace{2em}

This chapter assembles the evidence that motivates an AI-powered virtual
teaching assistant and positions this thesis against prior work. It proceeds
in two movements. \autoref{review:sec:workload} reviews three literatures that
have so far run in parallel: time-use surveys that quantify how much of an
educator's week non-teaching work consumes, econometric evidence that this
burden degrades teaching quality when it is uncompensated, and the
AI-in-education literature that has long nominated intelligent assistants as
the remedy without ever measuring the relief they deliver.
\autoref{review:subsec:gap} states the three-way gap at their intersection,
which this thesis addresses. \autoref{review:sec:rag} then reviews the
technical line of work --- retrieval-augmented generation and its refinements
--- that separates the system built here from the scripted chatbots surveyed
in the earlier education literature.

\section{Educator Workload, Administrative Burden, and AI-Assisted
         Intervention}
\label{review:sec:workload}

\subsection{Educator Time Allocation: The Measured Burden}
\label{review:subsec:burden}

The most comprehensive measurement of how educators spend their working week
is the OECD's Teaching and Learning International Survey (TALIS), whose 2024
cycle covered some 280\,000 educators across 55 education
systems~\cite{oecd25:Talis}. Across the OECD, full-time lower-secondary
teachers report a 41-hour working week, of which only 22.7 hours is classroom
teaching; general administrative work accounts for roughly 3 hours per week, a
share that has remained essentially stable since the 2018
cycle~\cite{oecd25:Talis}.

Singapore, the deployment context of this thesis, sits well above these
averages. Its teachers report a 47.3-hour working week --- about six hours
above the international average --- of which only 17.7 hours is teaching, with
8.2 hours going to lesson preparation, 6.4 hours to marking, and 4.0 hours to
administrative work~\cite{oecd25:TalisSG}. Teaching, in other words, occupies
a minority of the working week. The burden is also felt as such: 53\% of
Singapore teachers rank excessive administrative work as their foremost source
of occupational stress, with marking second at 49\%~\cite{oecd25:TalisSG}.
Internationally the administrative share varies widely --- from more than six
hours per week in Korea down to roughly an hour and a half in Finland and
France --- but it is persistent wherever it is high, having barely moved in
most systems between 2018 and 2024~\cite{oecd25:Talis}.
% TODO: verify the cross-system figures (Korea, Finland, France; optionally
% Japan ~5h, Australia 4.7h) against the published TALIS 2024 tables before
% submission, and cite the specific table numbers.

TALIS surveys school teachers, whereas this thesis deploys into university
courses; the two populations must be bridged explicitly rather than silently.
The higher-education evidence tells the same story at a larger scale: US
faculty in the National Study of Postsecondary Faculty (NSOPF) data report
working weeks on the order of 53 hours, with roughly a fifth going to
administrative and other non-instructional work, and economics faculty
specifically report about 6.2 hours per week of administration plus a further
3.7 hours of institutional service~\cite{wp07:Allgood, aer13:Allgood}. The
non-teaching burden is thus documented at both the school and university
levels; the Singapore school data establishes the local policy salience of the
problem, while the faculty studies carry the evidence for the university
setting in which this system operates.

\subsection{Non-Teaching Burden and Teaching Quality}
\label{review:subsec:quality}

Hours consumed would matter less if they were workload-neutral for students.
The central evidence that they are not comes from
\citet{apec15:GarciaGallego}~\cite{apec15:GarciaGallego}, who analyse a panel
of 604 professors at a medium-sized public Spanish university over 2002--2006
using random-effects GLS, relating teaching quality to research output and
administrative duties. Their key result is a distinction between two kinds of
administrative burden. Holding \emph{light} administrative positions ---
duties that carry no statutory teaching-load relief --- has a negative and
statistically significant effect on teaching quality ($-0.064$, $z = -2.47$),
whereas \emph{heavy} administrative positions, which come with a compensating
statutory load reduction, show no significant effect. The harm, in other
words, is determined not by the duty itself but by whether the time it
consumes is compensated. This is the pivotal finding for the present thesis:
a virtual teaching assistant that absorbs routine support load is precisely a
compensation mechanism, one that requires no statutory decree.

The same study clarifies a long-standing puzzle in the research--teaching
nexus literature. Earlier syntheses assuming a linear relationship found
little connection between research and teaching: \citet{rhe87:Feldman} reported a small
positive correlation~\cite{rhe87:Feldman}, \citet{rer96:Hattie} an association
near zero --- the relationship famously dismissed as an \enquote{enduring
myth}~\cite{rer96:Hattie} --- and \citet{jhe02:Marsh} a near-null
relation~\cite{jhe02:Marsh}. García-Gallego et al.\ show the relationship is
instead nonlinear: teaching quality rises with research output up to an
interior peak (an inverted U, with roughly two-thirds of the faculty below the
peak), and professors with no research at all are about five times more likely
than their colleagues to be among the worst-rated teachers (14.8\% versus
3.04\%)~\cite{apec15:GarciaGallego}. Under a linearity assumption, such a
relationship averages out to the earlier null results. Two accuracy caveats
apply when using this evidence. First, the study measures administrative
burden by positions held over years, not by hours, so it establishes a quality
effect but not a time mechanism. Second, its quality outcome is built from
student evaluations of teaching (SET), which are known to be susceptible to
expected-grade bias and to measure perceived rather than actual learning ---
a limitation that motivates the richer, behaviour-based telemetry proposed in
this thesis (\autoref{eval:sec:metrics}).

Where does the time for uncompensated duties come from? Allgood and Walstad's
time-allocation evidence shows that research and teaching do trade off
directly --- roughly ten additional weekly hours of research correspond to
about three fewer hours of teaching --- but that administrative and
\enquote{other} time does \emph{not} significantly displace teaching
hours~\cite{wp07:Allgood, aer13:Allgood}. Faculty absorb the administrative
load as overtime, which is why the cost remains invisible in time-use data
even as the quality evidence above shows it is real. The field's own normative
remedy points the same way: the AAUP holds that heavy service assignments
impair teaching effectiveness and that the correct response is workload relief
estimated in hours~\cite{aaup00:Workload}. An assistant that deflects routine
queries is exactly such hour-denominated relief, delivered by a system rather
than by policy.

\subsection{AI as a Workload Intervention: The Recognised Thread and Its
            Limits}
\label{review:subsec:ai}

That AI systems could supply this relief is not a new idea.
\citet{sust22:Ahmad}~\cite{sust22:Ahmad} survey AI applications in education
and organise them at two levels --- \emph{administrative} (admissions,
counselling, library services) and \emph{academic} (assessment, feedback,
tutoring) --- a taxonomy this thesis adopts: the system built here operates at
the academic level, with administrative spillover through the deflection of
routine questions. Their core claim is that AI applications
\enquote{minimize the administrative tasks of a teacher to invest more in
teaching and guiding students}, with chatbots singled out for operating around
the clock and reducing the question-answering burden on staff. The larger
systematic-review literature reaches a similar position:
\citet{ijethe19:ZawackiRichter}~\cite{ijethe19:ZawackiRichter}, reviewing
research on AI in higher education, find the field concentrated on
profiling, assessment and intelligent tutoring, while pointedly asking
\enquote{where are the educators?} --- the teacher-workload perspective is
largely absent from the systems the field builds and evaluates.

The practitioner side corroborates the intervention hypothesis. In TALIS 2024,
teachers in systems with high administrative load disproportionately believe
AI can automate administrative tasks, and 52\% across the OECD believe AI can
adapt materials to student abilities~\cite{oecd25:Talis}. Singapore is a
natural deployment context on this evidence: 76\% of its teachers have already
been trained in the use of AI, against single-digit percentages in some
European systems~\cite{oecd25:TalisSG}.

Two limits of this thread must be stated plainly. First, its evidentiary
weight is low: Ahmad et al.'s survey is a narrative review with a managerial
scope and no quantitative testing, a limitation its authors declare themselves
--- the claims, they note, still need to be quantitatively checked to be
generalisable~\cite{sust22:Ahmad}. Second, its technology is dated: the
chatbot precedents it rests on are keyword- and string-similarity systems
answering university FAQs from hand-maintained scripts, such as the AIML-based
assistant of \citet{icacci17:Ranoliya}~\cite{icacci17:Ranoliya} --- pre-LLM
retrieval, reviewed further in \autoref{review:sec:rag}. The thread thus
supplies the framing and the hypothesis, but neither the mechanism nor the
measurement.

\subsection{The Gap This Thesis Addresses}
\label{review:subsec:gap}

Read together, the three literatures interlock but do not close. The
time-accounting literature (\autoref{review:subsec:burden}) quantifies the
hours that non-teaching work consumes, but says nothing about its effects.
The quality literature (\autoref{review:subsec:quality}) quantifies the
effect --- uncompensated burden degrades teaching --- but measures duties in
positions held over years rather than hours, so it contains no time
mechanism; and because faculty absorb the burden as overtime, the displacement
never surfaces in time-use data. The AI-intervention literature
(\autoref{review:subsec:ai}) proposes the remedy but, by its own admission,
has never quantified the workload reduction it promises, and the systems it
surveys predate retrieval-augmented large language models. No study connects
the three: nobody has deployed the proposed intervention and measured, in
hours, the relief it delivers.

This thesis addresses that gap with two contributions. First, it builds the
modern form of the proposed intervention: a retrieval-augmented virtual
teaching assistant over institutional (NUS Canvas) course materials
(Chapters~\ref{ch:design} and~\ref{ch:implementation}). Second --- and unlike
prior work --- it instruments the intervention, logging every deflected query,
to produce a direct estimate of the instructor-hours returned per course per
week: the hour-denominated measurement that all three literatures lack. The
metric and its results are developed in \autoref{ch:evaluation}
(\autoref{eval:sec:metrics}). \autoref{review:tab:findings} summarises the
sources reviewed in this section and the specific findings this thesis builds
on.

\begin{table}[htb]
  \centering
  \caption{Summary of the workload, teaching-quality, and AI-intervention
           evidence used in this chapter.}
  \label{review:tab:findings}
  \footnotesize
  \renewcommand{\arraystretch}{1.3}
  \begin{tabular}{L{0.21\textwidth} L{0.13\textwidth} L{0.14\textwidth} L{0.4\textwidth}}
    \hline
    \textbf{Source} & \textbf{Type} & \textbf{Population} & \textbf{Key finding used} \\
    \hline
    OECD TALIS 2024, Singapore country note~\cite{oecd25:TalisSG} &
      Survey &
      SG lower-sec. teachers &
      4\,h/week admin; 17.7 of 47.3\,h/week teaching; admin the top stressor
      (53\%) \\
    OECD TALIS 2024, international report~\cite{oecd25:Talis} &
      Survey &
      55 systems &
      OECD average 3\,h/week admin; stable since 2018; AI seen by teachers as
      an admin remedy \\
    García-Gallego et al.~\cite{apec15:GarciaGallego} &
      Panel econometrics &
      604 professors, Spain &
      Uncompensated admin duties reduce teaching quality ($-0.064$);
      compensated duties show no effect; research--quality is an inverted U \\
    Allgood \& Walstad~\cite{wp07:Allgood, aer13:Allgood} &
      Survey + regression &
      US faculty &
      $\sim$6.2\,h/week admin; admin time absorbed as overtime rather than
      displacing teaching hours \\
    Hattie \& Marsh~\cite{rer96:Hattie} &
      Meta-analysis &
      58 studies &
      Research--teaching relation $\approx 0$ under a linear assumption \\
    AAUP~\cite{aaup00:Workload} &
      Policy statement &
      --- &
      Remedy for service burden is hour-estimated workload relief \\
    Ahmad et al.~\cite{sust22:Ahmad} &
      Narrative review &
      --- &
      Two-level AI taxonomy (administrative/academic); chatbots reduce Q\&A
      burden; claims self-declared as unquantified \\
    \hline
  \end{tabular}
\end{table}

\section{From Scripted Chatbots to Retrieval-Augmented Generation}
\label{review:sec:rag}

The chatbot systems surveyed by the AI-in-education literature of
\autoref{review:subsec:ai} share a common architecture: student questions are
matched against hand-authored scripts by keywords or string similarity, and
the response is the stored answer of the best-matching
pattern~\cite{icacci17:Ranoliya, sust22:Ahmad}. Such systems can deflect
narrowly anticipated FAQs, but they are brittle to paraphrase, cannot answer
questions whose answers were never scripted, and shift the maintenance burden
back onto staff --- every syllabus change requires editing scripts, which is
itself administrative work.

Retrieval-augmented generation (RAG)~\cite{neurips20:Lewis} replaces both
halves of that architecture. Questions and course materials are embedded in a
shared dense vector space, so retrieval is semantic rather than lexical and
survives paraphrase; and the answer is synthesised by a language model
conditioned on the retrieved passages, so the system can compose answers from
material nobody scripted, while remaining grounded in the course corpus rather
than in the model's parametric memory. The retrieval quality of such a
pipeline rests on the embedding model; this thesis uses
BGE-M3~\cite{arxiv24:Chen}, a multilingual, multi-granularity embedder suited
to heterogeneous course materials.

Naive RAG, however, has a known failure mode: retrieved chunks are prepended
to the prompt in full even though each is usually only partially relevant,
and language models demonstrably struggle to use long contexts, with relevant
information degraded or ignored depending on its position in the
prompt~\cite{tacl24:Liu}. \citet{iclr24:Xu}~\cite{iclr24:Xu} address this
with \textsc{Recomp}: retrieved passages are compressed before generation,
and --- through \emph{selective augmentation} --- an empty context is returned
when nothing relevant survives compression, rather than forcing the model to
attend to noise. This matters doubly for the present system, whose generation
runs on a small locally hosted model for which prompt length is the dominant
latency cost; the extractive compression stage adapted from this line of work
is described in \autoref{impl:sec:compression} and evaluated in
\autoref{eval:sec:results:compression}.

Positioned against both literatures, the distinction of this thesis is not
architectural novelty for its own sake: it is that the RAG pipeline is
deployed as the workload intervention that \autoref{review:subsec:ai}
identified, and instrumented to yield the measurement that
\autoref{review:subsec:gap} identified as missing.

\section{Summary}
\label{review:sec:summary}

Educator time is finite and heavily taxed by non-instructional work ---
quantified by TALIS 2024, with Singapore above the OECD average and
administrative work its teachers' top stressor~\cite{oecd25:TalisSG}. When
that burden is uncompensated, teaching quality measurably
suffers~\cite{apec15:GarciaGallego}, and because faculty absorb the load as
overtime rather than reduced teaching hours~\cite{wp07:Allgood}, the cost
stays hidden in time-use data. The field's own remedy is hour-denominated
load relief~\cite{aaup00:Workload}, and the AI-in-education literature has
long nominated intelligent assistants as the mechanism~\cite{sust22:Ahmad} ---
but concedes its claims are qualitatively argued and quantitatively
unverified, resting on pre-LLM keyword chatbots. This thesis supplies both
the modern mechanism --- a RAG-based virtual teaching assistant over NUS
Canvas course materials with extractive retrieval
compression~\cite{neurips20:Lewis, iclr24:Xu} --- and the missing
measurement: per-query deflection telemetry that estimates instructor-hours
returned, closing the loop between the workload, teaching-quality, and
AI-intervention literatures.
